\section*{Chapter \ref{ch:hf:extensions-of-the-inventory-framework} --- Extensions of the Inventory Framework}

\begin{solution}{exo:hf:extensions-of-the-inventory-framework:1}
$c=\sqrt{0.5e/210}=0.0804$; $\delta^b(2)=0.667+2.5c=0.868$, $\delta^a(2)=0.667-1.5c=0.546$.
\end{solution}

\begin{solution}{exo:hf:extensions-of-the-inventory-framework:2}
$q_0=\mu/(2\phi)=1$ unit.
\end{solution}

\begin{solution}{exo:hf:extensions-of-the-inventory-framework:3}
$17^3=4\,913$ points for three instruments; $17^{10}\approx2.0\times10^{12}$ for ten: the grid is out of reach, hence quadratic approximations.
\end{solution}

\begin{solution}{exo:hf:extensions-of-the-inventory-framework:4}
$\delta^a+\delta^b=2/k+c\bigl((2q+1)-(2q-1)\bigr)/2=2/k+c$, free of $q$: $1.333+0.080=1.414$.
\end{solution}

\begin{solution}{exo:hf:extensions-of-the-inventory-framework:5}
With positive correlation, a long position in asset 1 offsets part of the variance of the short in asset 2, so buying asset 1 lowers the penalty
$\phi q^\top\Sigma q$. With negative correlation the offsetting position is a short in asset 1: the market maker would bid less and offer more
aggressively.
\end{solution}

\begin{solution}{exo:hf:extensions-of-the-inventory-framework:6}
The move revalues the whole remaining position, not only the unit traded. After a bid fill the position is $q+1$, after an ask fill $q-1$; averaged
over the positions the market maker holds, the loss per fill beyond the unit itself cancels to a quadratic term in $q$ that the value function
absorbs, and what remains is half the move per fill.
\end{solution}

\begin{solution}{exo:hf:extensions-of-the-inventory-framework:7}
At $q=0$: bid 0.543, ask 0.872 (keen to buy); at $q=2$: 0.708 on both sides, exactly the no-drift depths at $q=0$. The drift moves the whole
schedule by $q_0=2$.
\end{solution}

\begin{solution}{exo:hf:extensions-of-the-inventory-framework:8}
Three possible errors. Widening by the whole move instead of half of it (a permanent move revalues the position too); a time step coarse enough to
change the answer (the chapter's coarse step favoured no widening, the fine one half the move); and scoring on mean P\&L while the policy maximises
P\&L less the penalty. Widen by half the permanent move, check with a quarter of the time step, and score the objective.
\end{solution}

\begin{solution}{pb:hf:extensions-of-the-inventory-framework:1}
\begin{enumerate}
\item A hard limit $|q|\le\bar q$; it makes the state space finite, so the linear system has a largest eigenvalue that dominates as the horizon grows.
\item See \cref{prop:hf:extensions-of-the-inventory-framework:stationary} (Perron--Frobenius and the spectral decomposition of $e^{M(T-t)}$).
\item See \cref{prop:hf:extensions-of-the-inventory-framework:closed}.
\item When $\beta=\sqrt{\phi ke/A}$ is small; worst gaps for $|q|\le10$: 0.034, 0.70 and 5.0.
\item One inventory vector and the charge $\phi q^\top\Sigma q$ on the whole book.
\item The substitution $h=\tfrac1k\ln\omega$ linearises every instrument's term only if all share the same $k$.
\item 0.572, 0.700 and 0.828.
\item 1\,000 common paths of five units of time, mids correlated 0.8, six penalties; inventory risk is the time average of $q^\top\Sigma q$.
\item 21\% less (24.3 against 30.8 at a mean profit of 673.3); 10\% to 22\% across penalties.
\item P\&L variance also comes from the randomness of fills and spread capture, which quoting jointly does not change.
\item The drift inside the control problem, acting through the value of holding inventory.
\item See \cref{prop:hf:extensions-of-the-inventory-framework:target}.
\item 355.6 against 335.7 (5.9\% more), holding 1.97 units on average against a target of 2.
\item Less: the market maker can only hold up to its bound, and the reward $-\phi(q-q_0)^2$ is then maximised at the bound.
\item See \cref{prop:hf:extensions-of-the-inventory-framework:adverse}.
\item Profit 262.7, 267.2, 258.9; objective 253.1, 258.6, 251.2 for none, 0.15 and 0.3.
\item With a flat quote filled with probability about 0.24 per step, the discrete simulation departs from the continuous-time model the proposition
solves; with a quarter of the step it agrees.
\item Adverse selection: it is present in every real market (chapter 1's quoter lost its spread to it), and the adjustment is a constant per side.
\item Mark out one's own fills at the horizon where the curve settles (One Quant Book 7, chapter 23): the settled mark-out is $\xi$.
\item Each is a change to one matrix: its size (bounds, assets), its diagonal (penalties, drift) or its off-diagonal (adverse selection).
\end{enumerate}
\end{solution}

\begin{solution}{iq:hf:extensions-of-the-inventory-framework:1}
Being long the fund is being long the index: skew the future's quotes as if you were already long it, keener to sell and less keen to buy, by the
hedge ratio times your fund inventory.
\iqlookfor{thinking in one book, with a hedge ratio.}
\end{solution}

\begin{solution}{iq:hf:extensions-of-the-inventory-framework:2}
The solution is $e^{M(T-t)}z$; far from the horizon the largest eigenvalue dominates, and the depths depend only on ratios of the dominant eigenvector.
\iqlookfor{the spectral argument, or an ergodic-control intuition.}
\end{solution}

\begin{solution}{iq:hf:extensions-of-the-inventory-framework:3}
Whether the move is a cost on the unit (widen by 0.3) or a permanent move of the mid (widen by 0.15); measure the mark-out curve and where it settles,
and check how much of the move reverts.
\iqlookfor{the distinction and a measurement plan.}
\end{solution}

\begin{solution}{iq:hf:extensions-of-the-inventory-framework:4}
Precompute: solve offline, store the depths in a table indexed by inventory (and a few signal buckets), or use the closed form; refresh the table
when parameters change, off the hot path.
\iqlookfor{precomputation and a lookup on the hot path.}
\end{solution}

\begin{solution}{iq:hf:extensions-of-the-inventory-framework:5}
The correlation between positions: offsetting positions look risky and concentrated ones look diversified. Measure $q^\top\Sigma q$ (or a
factor-model version) on the whole book.
\iqlookfor{portfolio variance, not a sum of variances.}
\end{solution}

\begin{solution}{iq:hf:extensions-of-the-inventory-framework:6}
The running reward $\mu q-\phi q^2=-\phi(q-\mu/(2\phi))^2+\mu^2/(4\phi)$: the problem is the zero-drift one in the shifted variable $q-q_0$, whose
optimum centres the inventory at zero.
\iqlookfor{completing the square.}
\end{solution}
